% Generated draft: requires mathematical and source review.
\documentclass[11pt]{article}
\usepackage[margin=25mm]{geometry}
\usepackage{fontspec}
\setmainfont{Latin Modern Roman}
\usepackage{amsmath,amssymb,mathtools,unicode-math}
\setmathfont{Latin Modern Math}
\usepackage{enumitem,needspace,xurl,hyperref}
\setlength{\parindent}{0pt}
\setlength{\parskip}{5pt}
\setlength{\emergencystretch}{3em}
\newcommand{\N}{\mathbb{N}}
\newcommand{\Z}{\mathbb{Z}}
\newcommand{\Q}{\mathbb{Q}}
\newcommand{\R}{\mathbb{R}}
\newcommand{\C}{\mathbb{C}}
\newcommand{\F}{\mathbb{F}}
\newcommand{\A}{\mathbb{A}}
\newcommand{\PP}{\mathbb{P}}
\newcommand{\E}{\mathbb{E}}
\newcommand{\eps}{\varepsilon}
\newcommand{\dd}{\,\mathrm{d}}
\begin{document}
\section*{UnsolvedMath ID: 30005645}
\textbf{Source problem code:} catalogue entry
\subsection{Definitions and mathematical statement}
Let \(p>0\) be a prime. A pointed stable curve of type \((g,n)\) over an algebraically closed field \(k\) of characteristic \(p\) is a connected, projective nodal curve of arithmetic genus \(g\), equipped with \(n\) ordered, distinct smooth marked points, satisfying stability: every irreducible component of its normalization, of genus \(h\) and with \(r\) special points coming from markings and branches of nodes, satisfies \(2h-2+r>0\). Thus the usual stable range is \(g,n\in\Z_{\geq 0}\) with \(2g-2+n>0\); this range is implicit in the terminology rather than explicitly specified in the supplied record.

The record denotes by \(\mathcal M_{g,n}^{\mathrm{adm}}\) a space of Frobenius-equivalence classes of such curves and by \(\Pi_{g,n}\) a space of isomorphism classes of their admissible fundamental groups. Writing \([X]\) for the intended curve class and \([\pi_1^{\mathrm{adm}}(X)]\) for the intended group class, the record supplies a natural continuous surjection
\[
\pi_{g,n}^{\mathrm{adm}}:\mathcal M_{g,n}^{\mathrm{adm}}\twoheadrightarrow\Pi_{g,n},\qquad [X]\longmapsto[\pi_1^{\mathrm{adm}}(X)].
\]
Here the notation \(\pi_1^{\mathrm{adm}}(X)\) names the admissible fundamental group referred to in the record; its construction is not supplied.

\textbf{Homeomorphism conjecture, as supplied.} For every positive characteristic \(p\) and every stable type \((g,n)\), the map \(\pi_{g,n}^{\mathrm{adm}}\), with the intended source and target topologies, is a homeomorphism. Equivalently, the supplied continuous surjection is bijective and its inverse
\[
(\pi_{g,n}^{\mathrm{adm}})^{-1}:\Pi_{g,n}\longrightarrow\mathcal M_{g,n}^{\mathrm{adm}}
\]
is continuous. The source distinguishes this from the weak Isom-version conjecture, which asserts only bijectivity of the underlying map of sets.

\textbf{Unresolved formulation qualifications.} The supplied record does not define Frobenius equivalence, the admissible fundamental group, either topology, or the treatment of curves over different algebraically closed fields. It also does not specify whether group isomorphisms must be continuous or preserve additional structure. These conventions are essential to a fully specified conjecture and must be recovered from the cited sources; the assertion above retains their intended meaning without selecting replacements. The excerpt presents the conjecture as relating the geometry of pointed stable curves and their moduli spaces to anabelian properties in positive characteristic; its subsequent local-moduli discussion is truncated and supplies no further mathematical assertion here.
\subsection{Short English statement}
For pointed stable curves in characteristic \(p>0\), the conjecture asks whether assigning the admissible fundamental group identifies the intended space of Frobenius-equivalence classes of curves with the intended space of group isomorphism classes by a homeomorphism. This includes both bijectivity and continuity of the inverse. Essential equivalence relations, group conventions, and topologies are missing from the supplied record, so this is a statement-only draft requiring mathematical review.
\subsection{Sources}
\begin{itemize}
\item Supplied local UnsolvedMath registry record, permanent ID \(30005645\), title \emph{Homeomorphism Conjecture for Moduli of Fundamental Groups}, exact source problem code \texttt{catalogue entry}, in the Oberwolfach Reports Open Problems collection. No individual proposer is identified in the record.
\item Supplied report citation: \emph{MFO--RIMS Tandem Workshop: Arithmetic Homotopy and Galois Theory (2024)}, Oberwolfach Reports, pp.~2377--2488, DOI \texttt{10.4171/owr/2023/42}. \url{https://doi.org/10.4171/owr/2023/42}. The record gives proposed year 2023, while its report citation displays 2024; this discrepancy is retained unresolved.
\item Supplied literature reference: Yang, ``moduli spaces of fundamental groups.'' \url{https://arxiv.org/abs/2010.01806}. The description is reproduced as supplied and has not been independently checked.
\item Supplied literature reference: Hu--Yang--Zong, ``topological structures and anabelian geometry.'' \url{https://arxiv.org/abs/2301.04864}. The description is reproduced as supplied and has not been independently checked.
\item Literature assessment dated exactly \texttt{2026-08-22}: the supplied curation labels the problem partially solved and reports substantial positive results in special loci and low-dimensional cases, while describing the assertion for all pointed stable curves in positive characteristic as conjectural. This is supplied dated curation derived from the OWR collection, not an independently verified literature assessment; the particular cases and hypotheses are not specified in the record.
\item The registry labels the statement \texttt{corrected\_verified} and reports that it was checked against the official Oberwolfach report. This is supplied source-verification metadata, not a verification performed for this draft.
\end{itemize}
\end{document}
